⚠️ 2) Problemas principais (que prejudicam o texto)
(A) O texto está longo demais e perde foco no meio
Você começa muito bem, com argumentação coesa, mas a partir de Joseph, Baboota, Herold… o texto vira praticamente “resumo de artigo por artigo”, sem integração temática.
O leitor sente isso como sequência de blocos soltos.
Falta:
síntese,
agrupamento temático,
análise crítica,
ligação de um estudo com o outro.
Isso é obrigatório numa revisão de literatura.

(B) Muitos parágrafos seguem a mesma estrutura
“Os autores fizeram X. Eles usaram Y. Obtiveram Z."
Isso torna a leitura monótona e pouco analítica.


🔥 3) Ajustes estratégicos que vão elevar seu texto a nível de dissertação
(1) Dê um propósito claro à revisão
Logo no início, você deve amarrar:
onde seu trabalho se encaixa nessa literatura?
Não faça isso no fim — faça no começo e no fim.

(2) Agrupe estudos por categorias, não por ordem cronológica
Isso melhora fluidez e reduz repetição.
Exemplo sugerido:
Revisões amplas — Horvat, Rico
Previsão de resultados — Joseph, Baboota, Awadallah
Análise de desempenho e tática — Herold, MLSA
Lesões — Majumdar
Deep learning e avanços recentes
Contexto brasileiro — SALab/FAME
Assim seu capítulo parece estruturado, não uma sequência de resumos.

(3) Diga por que cada estudo importa
Não apenas “o que fizeram”.
(4) Faça frases de transição
Elas são essenciais para parecer um texto único.
Agora você pula sem aviso de Horvat → Rico → Joseph → Baboota.
Eu posso escrever essas transições para você, se quiser.

❗ 4) Problemas pontuais que precisam ser corrigidos agora
“grandes volume de dados” → grandes volumes
“os autores destacam que diferentes algoritmos têm sido empregados” → está vago e repetitivo (Horvat fala isso sim, mas você poderia especificar o porquê).
Uso repetido de “Ainda assim,” em 3 parágrafos seguidos → varie.
O parágrafo do Rico está bom, mas ele está muito próximo do estilo do Horvat — sugiro reforçar o contraste entre eles.

🧠 Resumo da crítica honesta
→ O texto é tecnicamente correto e bem escrito, mas não é ainda o nível ideal de uma dissertação, porque falta síntese, estrutura temática e análise crítica.
Se você me disser “reescreva mantendo tudo que eu falei, mas com estrutura acadêmica de revisão de literatura de mestrado”, eu ajusto:
coerência,
categorização,
transitions,
redução de redundâncias,
clareza,
força argumentativa.
Posso inclusive:
reorganizar tudo em tópicos temáticos
manter cada citação
tornar seu texto muito mais elegante e sólido
